\documentclass{report}
\usepackage{amsfonts, amsmath}
\newcommand\R{{\mathbb R}}
\newcommand\Q{{\mathbb Q}}
\newcommand\N{{\mathbb N}}
\newcommand\ve{\varepsilon}
\newcommand\Co{{\mathcal C}}
\pagestyle{empty}
\begin{document}
\large
\centerline{\Huge Fisica Matematica II}
\renewcommand{\thefootnote}{ } 
\footnote{\sl Avete 2 ore di tempo. Ogni esercizio
vale dieci punti. La sufficienza si ottiene con un punteggio $\geq$
18. Solo le risposte chiaramente giustificate saranno prese in
considerazione.}
\renewcommand{\thefootnote}{\arabic{footnote}} 
\vskip.1cm
\centerline{\Large Seconda Prova di Autovalutazione}
\vskip1cm
\begin{enumerate}
\item Trovare la soluzione per
\[
\begin{split}
&u_{xx}+xe^{-t}=u_t\quad x\in[0,1], t>0\\
&u(x,0)=x \quad\quad x\in[0,1]\\
&u(t,0)=u(t,1)=0\quad\quad x\in[0,1].
\end{split}
\]
\item Trovare la soluzione per
\[
\begin{split}
&u_{tt}=u_{xx}+xe^{-t}\quad x\in\R,0<t\\
&u(x,0)=x;\;u_t(x,0)=0 \quad x\in \R.
\end{split}
\]
\item Si calcoli la trasformata di Fourier di
\[
e^{-ax^2}
\]
per $a>0$.
\item Si consideri l'equazione
\[
\begin{split}
&u_{xx}+e^{-x^2}e^{-t}\sin x=u_t\quad x\in\R,t>0\\
&u(x,0)=0\quad x\in\R.
\end{split}
\]
Si dimostri che $\sup_{(x,t)\in\R\times \R_+}|u(x,t)|\leq 1$.
\end{enumerate}
\end{document}
